%! Polynomial Functions and Complex Zeros — Unit 2, Lesson 2.4 — Group Activity (Answer Key)
\documentclass[10pt]{article}
\usepackage{precalculus-article}
\usepackage{precalculus-key}   % pulls in -boxes; do NOT also load -boxes
\newcommand{\TallMath}[1]{$\displaystyle #1\rule[-1.4em]{0pt}{3.2em}$}

\begin{document}

\pageheader{Unit 2, Lesson 2.4}{Group Activity --- Answer Key}

\vspace{0.12in}

\begin{scenariobox}[The Missing Zeros Bureau]{plumlight}
\small
Your group works the missing-zeros desk: every polynomial that comes in must
leave with \textit{all} of its zeros accounted for, whether or not the graph is
willing to show them. Tier E uses the pre-drawn graph of
$p(x) = x^3 - 3x^2 + 4x - 2$ below --- no tier draws anything.

\begin{center}
\begin{tikzpicture}[x=1.0cm, y=0.32cm]
  \draw[linegray, very thin, xstep=1, ystep=2] (-1,-4) grid (3,4);
  \draw[->, thick] (-1.3,0) -- (3.3,0) node[below right] {\small $x$};
  \draw[->, thick] (0,-4.8) -- (0,4.8) node[above] {\small $p(x)$};
  \foreach \x in {-1,1,2,3} \node[below] at (\x,-0.2) {\small \x};
  \foreach \y in {-4,-2,2,4} \node[left] at (0,\y) {\small \y};
  \draw[very thick, plum] plot [smooth] coordinates
    {(-0.4,-4.14) (-0.2,-2.93) (0,-2) (0.4,-0.82) (0.7,-0.33)
     (1,0) (1.3,0.33) (1.6,0.82) (2,2) (2.2,2.93) (2.4,4.14)};
  \fill[plum] (1,0) circle (2.5pt);
\end{tikzpicture}
\end{center}
\end{scenariobox}

\vspace{0.1in}

\boxguard
\begin{tcolorbox}[colback=white, colframe=black!40,
    title=\textbf{Tier R --- Remediate},
    fonttitle=\bfseries, arc=2mm, left=3mm, right=3mm, top=2mm, bottom=2mm]

\begin{enumerate}[itemsep=8pt]
  \item A polynomial function has degree 7. Counting multiplicities, how many
        complex zeros does it have? \quad \ans{7}

  \item Every case below uses the same move:
        degree $-$ (\# real zeros) $=$ \# non-real zeros. Fill the last column.

        \smallskip
        \renewcommand{\arraystretch}{1.5}
        \begin{center}
        \begin{tabular}{|c|c|c|}
        \hline
        \textbf{Degree} & \textbf{\# real zeros (with mult.)} & \textbf{\# non-real zeros} \\
        \hline
        4 & 4 & \ans{0} \\
        \hline
        4 & 2 & \ans{2} \\
        \hline
        5 & 3 & \ans{2} \\
        \hline
        6 & 0 & \ans{6} \\
        \hline
        \end{tabular}
        \end{center}

  \item Name the complex conjugate of each number. Only the sign in front of
        the $i$ changes.

        \begin{enumerate}[label=(\alph*), itemsep=4pt]
          \item $2 + 5i$ \quad \ans{$2 - 5i$}
          \item $-1 + i$ \quad \ans{$-1 - i$}
          \item $4i$ \quad \ans{$-4i$}
          \item $3 - 7i$ \quad \ans{$3 + 7i$}
        \end{enumerate}

  \item A polynomial with real coefficients has $6 - i$ as one of its zeros.
        Name another zero you get for free. \quad \ans{$6 + i$}
\end{enumerate}
\end{tcolorbox}

\vspace{0.1in}

\boxguard
\begin{tcolorbox}[colback=white, colframe=black!40,
    title=\textbf{Tier A --- Approaching Proficiency},
    fonttitle=\bfseries, arc=2mm, left=3mm, right=3mm, top=2mm, bottom=2mm]

\begin{enumerate}[itemsep=8pt]
  \item A polynomial $p$ has real coefficients, degree 3, leading coefficient
        1, and zeros $x = -1$ and $x = 2i$.

        \begin{enumerate}[label=(\alph*), itemsep=4pt]
          \item The third zero must be \ans{$-2i$}, because
                \ans{non-real zeros come in conjugate pairs}.
          \item Write $p(x)$ in expanded form.

                \begin{work}
                  p(x) &= (x+1)(x - 2i)(x + 2i) \\
                       &= (x+1)\big(x^2 - 4i^2\big) \\
                       &= (x+1)(x^2 + 4) \\
                       &= x^3 + x^2 + 4x + 4
                \end{work}

                $p(x) = $ \ans{$x^3 + x^2 + 4x + 4$}
        \end{enumerate}

  \item Classify each function as \textbf{even}, \textbf{odd}, or
        \textbf{neither}. Look at the exponents --- and remember a constant
        term rides on $x^0$.

        \begin{enumerate}[label=(\alph*), itemsep=4pt]
          \item $f(x) = 2x^6 - x^2 + 5$ \quad \ans{even}
          \item $g(x) = -x^7 + 3x$ \quad \ans{odd}
          \item $h(x) = x^4 + x^3$ \quad \ans{neither}
          \item $k(x) = x^2 - 9$ \quad \ans{even}
        \end{enumerate}

  \item Prove your answer to 2(b) by substituting $-x$ --- do not argue from
        the exponents this time.

        \begin{work}
          g(-x) &= -(-x)^7 + 3(-x) \\
                &= x^7 - 3x \\
                &= -\big(-x^7 + 3x\big) \\
                &= -g(x)
        \end{work}

        Therefore $g$ is \ans{odd}.
\end{enumerate}
\end{tcolorbox}

\vspace{0.1in}

\boxguard
\begin{tcolorbox}[colback=white, colframe=black!40,
    title=\textbf{Tier E --- Extension},
    fonttitle=\bfseries, arc=2mm, left=3mm, right=3mm, top=2mm, bottom=2mm]

\begin{enumerate}[itemsep=8pt]
  \item Read the graph of $p(x) = x^3 - 3x^2 + 4x - 2$ at the top of the page,
        then reconcile it with the algebra.

        \begin{enumerate}[label=(\alph*), itemsep=4pt]
          \item Number of $x$-intercepts on the graph: \ans{1}
          \item So the number of \textbf{real} zeros is \ans{1}
          \item The degree is \ans{3}, so the total number of complex
                zeros is \ans{3}
          \item Therefore $p$ has \ans{2} non-real zeros.
        \end{enumerate}

  \item Confirm the visible zero, then find the hidden pair. It happens that
        $p(x) = (x - 1)(x^2 - 2x + 2)$.

        \begin{work}
          p(1) &= (1)^3 - 3(1)^2 + 4(1) - 2 \\
               &= 1 - 3 + 4 - 2 \\
               &= 0
        \end{work}

        \begin{work}
          x^2 - 2x + 2 &= 0 \\
                     x &= \frac{2 \pm \sqrt{4 - 8}}{2} \\
                       &= \frac{2 \pm 2i}{2} \\
                       &= 1 \pm i
        \end{work}

        The two non-real zeros are \ans{$1 + i$ and $1 - i$}. Are they a conjugate pair?
        \quad \ans{Yes} / \textbf{No}

  \item Explain why \textit{every} polynomial of odd degree with real
        coefficients must cross the $x$-axis at least once. Use the pairing
        rule in your answer.

        \vspace{0.05in}
        \ansline{An odd-degree polynomial has an odd total number of zeros. Non-real zeros}
        \ansline{always come in conjugate pairs, so the non-real count is even --- an even}
        \ansline{count can never use up an odd total, so at least one zero must be real.}
\end{enumerate}
\end{tcolorbox}

\end{document}
